\begin{axiombox}{Axiom (E1, Reflexivity of Equality)}
\begin{axiom}[Reflexivity of Equality]
\label{ax:equality-reflexivity}
\[
\forall x,\ x=x.
\]
\end{axiom}
\end{axiombox}

\begin{remark*}[Standard quantified statement]
\[
\operatorname{Reflexive}(=)
\Longleftrightarrow
\forall x,\ x=x.
\]
\end{remark*}

\begin{remark*}[Interpretation]
Every object is itself. This is the primitive equality axiom used before
symmetry, transitivity, or arithmetic structure is derived.
\end{remark*}

\NoLocalDependencies

\begin{theorembox}{Theorem (E2, Symmetry of Equality)}
\begin{theorem}[Symmetry of Equality]
\label{thm:equality-symmetry}
\[
\forall x\,\forall y,\ (x=y \Longrightarrow y=x).
\]
\end{theorem}
\end{theorembox}

\begin{remark*}[Standard quantified statement]
\[
\operatorname{Symmetric}(=)
\Longleftrightarrow
\forall x\,\forall y,\ (x=y \Longrightarrow y=x).
\]
\end{remark*}

\begin{remark*}[Interpretation]
If $x$ is the same object as $y$, then $y$ is the same object as $x$.
Symmetry follows from Leibniz's law by substituting into the property
$P(t) \equiv (t=x)$ and using reflexivity.
\end{remark*}

\begin{dependencies}
\begin{itemize}
  \item \hyperref[ax:leibniz-law]{Leibniz's Law}
  \item \hyperref[ax:equality-reflexivity]{Reflexivity of Equality}
\end{itemize}
\end{dependencies}

\begin{theorembox}{Theorem (E3, Transitivity of Equality)}
\begin{theorem}[Transitivity of Equality]
\label{thm:equality-transitivity}
\[
\forall x\,\forall y\,\forall z,\,
\bigl(x=y \land y=z \Longrightarrow x=z\bigr).
\]
\end{theorem}
\end{theorembox}

\begin{remark*}[Standard quantified statement]
\[
\operatorname{Transitive}(=)
\Longleftrightarrow
\forall x\,\forall y\,\forall z,\,
\bigl(x=y \land y=z \Longrightarrow x=z\bigr).
\]
\end{remark*}

\begin{remark*}[Interpretation]
Sameness chains: if the first object is the second and the second is the
third, then the first is the third.
\end{remark*}

\begin{dependencies}
\begin{itemize}
  \item \hyperref[ax:leibniz-law]{Leibniz's Law}
  \item \hyperref[ax:equality-reflexivity]{Reflexivity of Equality}
  \item \hyperref[thm:equality-symmetry]{Symmetry of Equality}
\end{itemize}
\end{dependencies}

\begin{corollarybox}{Corollary (Equality Is an Equivalence Relation)}
\begin{corollary}[Equality Is an Equivalence Relation]
\label{cor:equality-is-equivalence-relation}
The equality relation is reflexive, symmetric, and transitive.
\end{corollary}
\end{corollarybox}

\begin{remark*}[Standard quantified statement]
\[
\operatorname{EquivalenceRelation}(=)
\Longleftrightarrow
\operatorname{Reflexive}(=)\land\operatorname{Symmetric}(=)\land\operatorname{Transitive}(=).
\]
\end{remark*}

\begin{remark*}[Interpretation]
Equality is the strictest equivalence relation: it groups only objects that
are literally the same.
\end{remark*}

\begin{dependencies}
\begin{itemize}
  \item \hyperref[ax:equality-reflexivity]{Reflexivity of Equality}
  \item \hyperref[thm:equality-symmetry]{Symmetry of Equality}
  \item \hyperref[thm:equality-transitivity]{Transitivity of Equality}
\end{itemize}
\end{dependencies}
